\documentclass[11pt]{article}

\usepackage[margin=1.1in]{geometry}
\usepackage{amsmath,amssymb,amsthm,mathrsfs}
\usepackage[colorlinks=true,linkcolor=black,urlcolor=black,citecolor=black]{hyperref}

\newtheorem{theorem}{Theorem}
\newtheorem{lemma}[theorem]{Lemma}
\newtheorem{proposition}[theorem]{Proposition}
\newtheorem{corollary}[theorem]{Corollary}
\theoremstyle{definition}
\newtheorem{definition}[theorem]{Definition}
\theoremstyle{remark}
\newtheorem{remark}[theorem]{Remark}

\newcommand{\C}{\mathbf C}
\newcommand{\Q}{\mathbf Q}
\newcommand{\Z}{\mathbf Z}
\newcommand{\PP}{\mathbf P}

\title{Frame transport for framed formal monodromy:\\
  the lemma, its hypotheses, and the receiver problem\\
  \large Third version, reporting the source checks and correcting the second}
\author{Memo for a reader of \emph{Irrationality of Cubic Threefolds after One
    Stabilization}}
\date{15 August 2026}

\begin{document}
\maketitle

\begin{abstract}
Proposition~4.7 of the draft transports framed formal monodromy across the
comparison isomorphisms of Iritani and Iritani--Koto, and the transport step
itself is asserted in one sentence. This memo states and proves the missing
lemma over a differential field, with its hypotheses made explicit; shows that
those hypotheses are jointly unsatisfiable over the coefficient receiver the
draft builds, so that the lemma as proved has no verified instance in the
application; corrects the receiver; and reduces the remaining question to one
inequality that depends on a normalization the draft is free to choose. The
inequality fails under the draft's current instruction to take the separating
weight parameter large, and holds at the cubic endpoint under the smallest
admissible choice. Section~\ref{sec:iterated} examines an iterated receiver that
would satisfy the inequality trivially, checks it against the sources, and finds
that it needs a semipositivity hypothesis on the ambient variety which weak
factorization does not supply.

The second version of this memo proposed the framing theory of Hinault, Yu,
Zhang and Zhang as the probable repair, on the ground that its decomposition
gauge is regular in the loop coordinate. Three compatibility checks against the
pinned sources have since been run, and the route does not survive them;
Section~\ref{sec:hyzz} now reports what those checks found. Two consequences are
worth more than the route was. First, no receiver is needed for the comparison
isomorphisms at all: Iritani's $\Psi$ and Iritani--Koto's $\Phi$ are polynomial
in the loop coordinate with inverses in the same ring, so the transport theorem
below applies to them with $\Gamma=0$. Second, the diagnosis of the second
version was wrong --- the pro-Laurent bulk gauge is not the draft's own choice
but Iritani--Koto's own fundamental solution, which the sources meet and handle
by Birkhoff factorization. What is left is a single statement about one bulk
displacement, recorded in Section~\ref{sec:hyzz}.
Section~\ref{sec:corrections} lists what the earlier versions
got wrong; readers of those versions should start there.
\end{abstract}

\section{The assertion at issue}

References are to Section~4 of the draft, \emph{Framed formal monodromy}. I name
statements by their semantic labels as well as by number, because the edits
proposed in Section~\ref{sec:placement} renumber the section: Definition~4.1
(\texttt{def:framed-sixth-multiplicity}), Definition~4.2
(\texttt{def:pro-laurent-gauge-group}), Lemma~4.3
(\texttt{lem:pv-base-change}), Lemma~4.5 (\texttt{lem:formal-base-shift}),
Proposition~4.7 (\texttt{prop:framed-operations}), Lemma~4.9
(\texttt{lem:divisor-tagging}).

Lemma~4.3 ends with the sentence
\begin{quote}
  \emph{If the gauge uses only integral powers of $z$ at every finite level,
  this conjugacy preserves the original-$z$-disc frame,}
\end{quote}
whose proof is the single sentence ``the deck action fixes every integral power
of $z$ at each finite level, so the compatible gauge is single-valued for the
framed turn.'' Lemma~4.5 then \emph{defines}
\[
  M^{\mathrm{bulk}}=G\,M^{\mathrm{small,shifted}}_{\mathrm f,V}G^{-1}
  \tag{4.1c}
\]
and Proposition~4.7 concludes that ``each center target summand has the framed
monodromy of its specialized small connection.'' The conclusion does not follow
from the display: conjugation preserves a characteristic polynomial for trivial
reasons, but says nothing about whether the conjugated matrix is the intrinsic
turn operator on horizontal solutions of the gauged module. What is missing is a
transport lemma together with an algebra in which the turn acts on solutions.

\section{The lemma over a differential field}
\label{sec:lemma}

Notation for $\operatorname{Exp}_V$, $K=\C\oplus V$ and $\Omega_V$ is the
draft's, from (4.0) and (4.0a). Recall that $\operatorname{Exp}_V$ is defined on
$(K,+)$, not on $\Omega_V$.

\subsection*{The ground field}

Let $\Gamma$ be a totally ordered abelian group and let
\[
  \mathcal H=\Omega_V(\!(\Gamma\oplus\Z e_z)\!)
\]
be the Hahn field of series with well-ordered support and coefficients in
$\Omega_V$, the value group ordered lexicographically with the $e_z$-coordinate
\emph{last}, as in the draft. Write $z=x^{e_z}$ and let $\theta$ multiply the
coefficient at $(\gamma,k)$ by $k$; this is a derivation because $k$ is additive
on the value group, and $\theta z=z$. Its constants are
\[
  \mathcal C:=\mathcal H^{\theta=0}=\Omega_V(\!(\Gamma)\!).
\]
Taking $\Gamma=0$ gives $\mathcal H=\Omega_V(\!(z)\!)$. For $e\ge1$ put
$\mathcal H_e=\Omega_V(\!(\Gamma\oplus\frac1e\Z e_z)\!)$, a degree-$e$ extension
containing $w=x^{e_z/e}$, and $\theta_w=e\theta$, so $\theta_ww=w$. All residues
below are taken with respect to $\theta_w$ on the ramified side; the draft's
$M_{\mathrm{RS},V}$ uses the same convention, and the two differ by the factor
$e$ if one is careless.

\subsection*{The solution algebra and the turn}

\begin{definition}
\label{def:solution-algebra}
Assume the exponential factors under consideration have \emph{constant}
coefficients, that is $\varphi\in w^{-1}\Omega_V[w^{-1}]$. The \emph{universal
formal solution algebra} $\mathcal U_e$ is the $\mathcal H_e$-algebra generated
by symbols
\[
  w^\rho\ (\rho\in K),\qquad \ell,\qquad E(\varphi),
\]
subject to $w^\rho w^{\rho'}=w^{\rho+\rho'}$, $w^n=$ the element $w^n$ of
$\mathcal H_e$ for $n\in\Z$, $E(\varphi)E(\varphi')=E(\varphi+\varphi')$,
$E(0)=1$, with
\[
  \theta_w(w^\rho)=\rho\,w^\rho,\qquad
  \theta_w(\ell)=1,\qquad
  \theta_w\bigl(E(\varphi)\bigr)=\theta_w(\varphi)E(\varphi).
\]
\end{definition}

The two hypotheses buried in that definition --- constant coefficients in
$\varphi$, and exponents $\rho$ in $K$ rather than in $\Omega_V$ --- are not
bookkeeping. The first is false for the modules the draft transports and is
discussed at Remark~\ref{rem:constants-fails}; the second is forced because
$\operatorname{Exp}_V$ is defined only on $K$, so a module whose residues leave
$K$ needs $\operatorname{Exp}$ rebuilt over a larger field.

\begin{lemma}[Structure]
\label{lem:structure}
Fix representatives $\Lambda\subset K$ for $K/\Z$ with $0\in\Lambda$. Then
$\mathcal U_e$ is free as an $\mathcal H_e$-module with basis
$\{w^\rho\ell^mE(\varphi)\}$, $\rho\in\Lambda$, $m\ge0$. In particular
$\mathcal U_e\ne0$ and $\mathcal H\subset\mathcal U_e$.
\end{lemma}

\begin{proof}
One cannot argue that $K/\Z$ is torsion-free; it is not, since it contains
$\Q/\Z$, and the rational classes are exactly the ones $\nu_6$ counts. Instead
use base change. The group algebra $\mathcal H_e[K]$ is free over
$\mathcal H_e[\Z]=\mathcal H_e[t,t^{-1}]$ on any set $\Lambda$ of coset
representatives, and $\mathcal U_e$ is obtained from it by the base change
$\mathcal H_e[t,t^{-1}]\to\mathcal H_e$, $t\mapsto w$, a surjection with kernel
$(t-w)$ generated by a nonzerodivisor. Base change of a free module is free on
the image of the same basis. The crossed multiplication $w^\rho
w^{\rho'}=w^nw^{\rho''}$ is well defined, and the derivation is consistent with
the integer identification because $w^n$ sits at $e_z$-exponent $n/e$ and
$\theta_w$ multiplies it by $n$. The factors $E(\varphi)$ form the group algebra
of a torsion-free group and $\ell$ is a polynomial variable.
\end{proof}

\begin{lemma}[The turn]
\label{lem:turn}
There is a unique automorphism $\sigma$ of $\mathcal U_e$ with
\[
  \sigma(w)=e^{2\pi i/e}w,\quad
  \sigma(w^\rho)=\operatorname{Exp}_V(\rho/e)\,w^\rho,\quad
  \sigma(\ell)=\ell+\tfrac{2\pi i}{e},\quad
  \sigma\bigl(E(\varphi)\bigr)=E(\sigma\varphi),
\]
acting on $\mathcal H_e$ by multiplying the coefficient at $(\gamma,j/e)$ by
$e^{2\pi ij/e}$. It fixes $\mathcal H$, hence $\mathcal C$, pointwise, and
commutes with $\theta_w$.
\end{lemma}

\begin{proof}
The action on $\mathcal H_e$ multiplies coefficients by a character of the value
group, hence is a support-preserving ring automorphism, and is the identity on
$\mathcal H$ because there $j\in e\Z$. Multiplicativity in $\rho$ is the
homomorphism property of $\operatorname{Exp}_V$ on $(K,+)$, and the two
prescriptions agree on integral exponents because
$\operatorname{Exp}_V(n/e)=e^{2\pi in/e}$ by (4.0a). Commutation with
$\theta_w$ is not a diagonality argument, since neither
$\sigma(\ell)=\ell+2\pi i/e$ nor $\theta_w(\ell^m)=m\ell^{m-1}$ is diagonal; it
is the direct check
$\sigma\theta_w(\ell^m)=m(\ell+\tfrac{2\pi i}e)^{m-1}=\theta_w\sigma(\ell^m)$,
together with the diagonal cases.
\end{proof}

\begin{lemma}[Constants]
\label{lem:constants}
Under the constant-coefficient hypothesis of
Definition~\ref{def:solution-algebra}, $\mathcal U_e^{\theta_w=0}=\mathcal C$.
\end{lemma}

\begin{proof}
Let $v$ be the Hahn valuation of $\mathcal H_e$, additive because the
coefficient field is a field; $\theta_w$ multiplies coefficients by integers and
never enlarges a support, so $v(\theta_wc)\ge v(c)$.

Write $x=\sum c_{\rho,m,\varphi}w^\rho\ell^mE(\varphi)$ in the basis of
Lemma~\ref{lem:structure} with $\theta_wx=0$; the equation splits over
$\varphi$. For $\varphi\ne0$, let $M$ be the top $\ell$-degree in that block;
for each $\rho$,
$\theta_w(c_{\rho,M})+\rho c_{\rho,M}+\theta_w(\varphi)c_{\rho,M}=0$. Every
monomial of $\theta_w(\varphi)$ has $\Gamma$-part $0$ and strictly negative
$e_z$-part, so $v(\theta_w(\varphi)c)<v(c)\le v(\theta_wc+\rho c)$ for $c\ne0$
and the lowest term cannot cancel; hence $c_{\rho,M}=0$, and descending in $M$
kills the block.

For $\varphi=0$ the coefficient equations are
$\theta_w(c_{\rho,m})+\rho c_{\rho,m}+(m+1)c_{\rho,m+1}=0$. Descending in $m$:
the top-degree coefficient satisfies the homogeneous equation
$\theta_w(c)+\rho c=0$, and on each monomial this reads $(j+\rho)c=0$ with
$j\in\Z$, so $c=0$ unless $\rho\in\Z$, that is unless $\rho=0$. For $\rho=0$ the
top coefficient is $\theta_w$-constant and the next equation gives
$Mc_{0,M}=-\theta_w(c_{0,M-1})$, whose right side has zero component in
$e_z$-degree $0$; comparing that component gives $c_{0,M}=0$ for $M\ge1$. What
survives is $\rho=0$, $m=0$, $\varphi=0$ with $\theta_w$-constant coefficient,
and the $\theta_w$-constants of $\mathcal H_e$ are those supported in
$e_z$-degree zero, namely $\mathcal C$.
\end{proof}

\begin{remark}[The constant-coefficient hypothesis is necessary, and fails in
the application]
\label{rem:constants-fails}
Allow $\varphi$ a Novikov coefficient of positive weight, $\varphi=x^\gamma
w^{-1}$ with $\gamma\in\Gamma$, $\gamma>0$ --- the shape of every exponential
factor the draft actually meets, since the eigenvalues $\lambda$ of $c_1\star$
have positive Novikov weight. Then $v(\theta_w\varphi)>0$ in the
coefficient-dominant order and the domination argument reverses. Worse, the
equation $\theta_w(c)=\theta_w(\varphi)c$ is solved by
\[
  c=\sum_{n\ge0}\frac{(-1)^n}{n!}\,x^{n\gamma}w^{-n},
\]
whose support $\{(n\gamma,-n/e)\}_{n\ge0}$ is increasing, hence well ordered, so
$c\in\mathcal H_e$ and $cE(\varphi)$ is a $\theta_w$-constant outside
$\mathcal C$. Lemma~\ref{lem:constants} is therefore false in that regime, and
with it Proposition~\ref{prop:framed-is-turn}. The structural reason is worth
stating on its own: \emph{in the draft's coefficient-dominant order,
$\exp(-\lambda/z)$ is already a unit of the coefficient field}, so adjoining
$E(\varphi)$ as a free symbol duplicates an existing element and inflates the
constants. In that order the exponential factors are not symbols at all, and the
entire irregularity of the formal solution sits in its divergent unit part,
which is exactly where the rate comparison of Section~\ref{sec:receiver} lives.
\end{remark}

\subsection*{Framed operators and transport}

\begin{proposition}[The framed operator is the matrix of the turn]
\label{prop:framed-is-turn}
Let $D$ be a differential module over $\mathcal H$, free of rank $n$, with
connection matrix $\Omega$, and let $Y\in\mathrm{GL}_n(\mathcal U_e)$ satisfy
$\theta Y=-\Omega Y$. Then
$\operatorname{Sol}(D)=\{v:\theta v+\Omega v=0\}=Y\mathcal C^{\,n}$, free of
rank $n$ and $\sigma$-stable; the matrix $M$ with $\sigma(Y)=YM$ lies in
$\mathrm{GL}_n(\mathcal C)$ and is well defined up to
$\mathrm{GL}_n(\mathcal C)$-conjugacy. For $D$ over $\Omega_V(\!(z)\!)$ in
Levelt--Turrittin form, $M$ is the draft's framed operator: in a solution basis
adapted to one deck orbit of length $d$ the matrix of $\sigma$ is block-cyclic
with the draft's intertwiners $A_i$ as blocks, and the cyclic-block determinant
gives $\chi(X)=\det(X^dI-U)$ with $U=A_{d-1}\cdots A_0$, which is (4.0b).
\end{proposition}

\begin{proof}
From $\theta(Y^{-1})=Y^{-1}\Omega$ one gets $\theta(Y^{-1}v)=Y^{-1}(\theta
v+\Omega v)$, so horizontality of $v$ means $Y^{-1}v$ has $\theta_w$-constant
entries, and Lemma~\ref{lem:constants} identifies those with $\mathcal C$. Since
$\sigma$ commutes with $\theta$ and fixes $\Omega$, it preserves
$\operatorname{Sol}(D)$; so $M=Y^{-1}\sigma(Y)$ has entries in $\mathcal C$ and
is invertible, being a product of invertibles, with inverse
$\sigma(Y^{-1})Y$. Two solution matrices differ on the right by an element of
$\mathrm{GL}_n(\mathcal C)$, which conjugates $M$.
\end{proof}

\begin{theorem}[Frame transport]
\label{thm:transport}
Let $D,D'$ be as in Proposition~\ref{prop:framed-is-turn}, with solution
matrices $Y,Y'$, and let $G\in\mathrm{GL}_n(\mathcal H)$ satisfy
$\Omega'=G\Omega G^{-1}-\theta(G)G^{-1}$. Then $GY$ is a solution matrix for
$D'$ and $\sigma(GY)=(GY)M$. Consequently
\[
  M_{\mathrm f,V}(D')=C_0^{-1}M_{\mathrm f,V}(D)\,C_0
  \qquad\text{for some }C_0\in\mathrm{GL}_n(\mathcal C),
\]
so the two framed operators have the same characteristic polynomial in
$\mathcal C[X]$ and the same multiplicity at every root of unity.
\end{theorem}

\begin{proof}
$\theta(GY)=\theta(G)Y-G\Omega Y=-\Omega'(GY)$. Since $\sigma$ fixes
$\mathcal H$ pointwise, $\sigma(G)=G$, so $\sigma(GY)=G\sigma(Y)=GYM$. Writing
$Y'=(GY)C_0$ with $C_0\in\mathrm{GL}_n(\mathcal C)$ gives $M'=C_0^{-1}MC_0$.
\end{proof}

\begin{remark}[Which object is $M_{\mathrm f,V}$]
The framed operator is an automorphism of the horizontal module; its matrix
depends on a choice of \emph{solution} frame and is well defined up to
conjugacy over the constants. The conjugate $GM G^{-1}$ appearing in the draft's
(4.1c) is not the matrix of the turn in any solution frame --- $Y'$ differs from
$GY$ by a constant matrix, not by $G$ --- but the matrix of the transported
automorphism read in the \emph{module} frame of $D'$. Both readings are
legitimate, and they give the same characteristic polynomial, but the draft
should say which it means, since conflating them is exactly the step at issue.
\end{remark}

\begin{remark}[The hypothesis on $G$ is the whole content]
Since $\sigma$ fixes $\mathcal H$ pointwise, $\sigma(G)=G$ is automatic once $G$
has entries in $\mathcal H$; what matters is that $G$ lies in $\mathcal H$ and
not in $\mathcal U_e$. No hypothesis-free statement exists: $Y$ itself is a
gauge over $\mathcal U_e$ carrying $D$ to the trivial module, whose framed
operator is the identity. The rank-one case shows what survives: if
$g\in\mathcal H^\times$ and $\theta(g)/g=\delta$, then each monomial gives
$(k-\delta)g_{(\gamma,k)}=0$, so $\delta\in\Z$; a gauge over $\mathcal H$ shifts
residues only by integers, which $\operatorname{Exp}_V$ kills.
\end{remark}

\subsection*{The scissors: these hypotheses have no common instance in the
application}
\label{sec:scissors}

Theorem~\ref{thm:transport} is true, and in the application it is empty. Its two
hypotheses pull in opposite directions.

\begin{itemize}
\item If $\Gamma=0$, so $\mathcal H=\Omega_V(\!(z)\!)$, then Levelt--Turrittin
  supplies solution matrices, but the pro-Laurent comparison gauge has
  unbounded negative $z$-order and is not in $\mathrm{GL}_n(\mathcal H)$.
\item If $\Gamma\ne0$ with the draft's coefficient-dominant order, the gauge
  lies in $\mathcal H$, but $\mathcal H$ is not a formal Laurent series field
  over its constants and Levelt--Turrittin is unavailable. Concretely, the
  $e_z$-component of the Hahn valuation is not a valuation on $\mathcal H$:
  for $f=x^{(0,0)}+x^{(1,-100)}$ and $g=-x^{(0,0)}$ its values on $f$ and $g$
  are $0$ while its value on $f+g$ is $-100$. So there is no Newton polygon, no
  slope filtration and no Turrittin algorithm. Nor can one descend to the
  bounded-$z$-order elements: $f=1-zx^{-\gamma}$ has inverse
  $-\sum_{k\ge1}z^{-k}x^{k\gamma}$, so they do not form a field.
\end{itemize}

An abstract Picard--Vessiot ring for a module over $\mathcal H$ does exist, with
constants $\mathcal C$ after passing to the divisible hull. So the missing
ingredient is \emph{not} a solution algebra. It is a canonical Levelt--Turrittin
normal form over $\mathcal H$, singling out which element of the differential
Galois group deserves to be called the turn. That is the precise shape of the
remaining problem, and it also shows why a purely valuation-theoretic argument
cannot supply it.

\section{The receiver}
\label{sec:receiver}

The draft's receiver, in the proof of Proposition~4.7, is
\[
  \mathscr R_j=\Omega_{V,j}\otimes_{H_{0,j}}H_{\mathrm{tot},j},
  \qquad
  H_{\mathrm{tot},j}=\C(\!(\Gamma^{\mathrm{coeff}}_j\oplus\Z e_z)\!).
  \tag{4.4c}
\]

\subsection*{Why the tensor product cannot simply be dissolved}

The map $a\otimes h\mapsto a\cdot h$ into the Hahn field with $\Omega_{V,j}$
coefficients is not well defined: the copy of $H_{0,j}$ inside $\Omega_{V,j}$
consists of degree-zero constants, the copy inside $H_{\mathrm{tot},j}$ of
monomial series, and the tensor product exists precisely to identify them. Nor
can the two copies be assigned different jobs, with module coefficients read as
constants and gauge bookkeeping as monomials: the gauge relation
$\Omega'=G\Omega G^{-1}-\theta(G)G^{-1}$ ties $G$ and $\Omega$ through the same
Novikov variables, so one copy is forced. With the constants reading the gauge
leaves the field; with the monomial reading the exponential factors become
$\Gamma$-supported and Remark~\ref{rem:constants-fails} applies.

Two things should also be said about $\mathscr R_j$ itself. It is a tensor
product of two field extensions and need not be a domain, and the draft reads
characteristic polynomials and algebraic multiplicities in $\mathscr R_j[X]$;
that deserves an explicit remark, which the draft partly gives by proving each
factor injects.

\subsection*{A better receiver}

Extend the $\Gamma$-valuation of $H_{0,j}$ to its algebraic closure and then to
$\Omega_{V,j}$, and take the field of series $\sum_kc_kz^k$ with
$c_k\in\Omega_{V,j}$ whose support $\{(v_\Omega(c_k),k)\}$ is well ordered in
the chosen order. This has a single copy of the coefficients, is a field rather
than a tensor product of two fields, hosts the pro-Laurent gauge, and hosts the
exponential factors as honest units. It is strictly better than (4.4c). It does
not by itself supply Levelt--Turrittin, and it does not remove the rate
criterion below.

\subsection*{The two rates}

Weight a monomial of coefficient weight $d$ and $z$-power $k$ by $a\,d+b\,k$,
$a,b>0$, and let $\varepsilon$ be the minimal weight of a generator of $J_j$.

\begin{itemize}
\item \emph{Pro-Laurent gauges gain one filtration unit per unit of $z$-pole.}
  By the recursion (4.1a) and the draft's own conclusion, $G_\alpha$ has no
  power of $z$ below $-|\alpha|$ and a term of bulk degree $m$ lies in $F^mB$;
  every monomial of $J_j^N$ has weight at least $N$. Their weights grow like
  $m(a-b)$, which together with the draft's finite-below argument gives well
  ordered support when $a>b$.
\item \emph{Splitting gauges lose $w(\Delta\lambda)$ per unit of $z$-power.}
  This is the draft's own construction (4.9d): at order $n$ the block
  off-diagonal coefficient is removed by a Sylvester equation $[K_0,A_n]=B_n$
  with $K_0$ the leading matrix, so one division by an eigenvalue difference per
  order, and $w(A_n)\ge-n\,w(\Delta\lambda)$. The draft's displayed cubic entry
  $-14/(81r^2)$ at $z$-order two exhibits the rate: one factor of $r^{-1}$ per
  unit of $z$-power. Three caveats. After a ramification $z=w^e$ the divisions
  occur once per $w$-order, so the loss is $e\,w(\Delta\lambda)$ per unit of
  $z$-power; at the cubic the draft proves the exponential blocks are
  unramified in $z$, so $e=1$ there. The claim that the regular-singular
  normalization costs nothing needs the homogeneity argument that residue
  eigenvalues are rational, verified in the cubic block where
  $\det L_s=(s+1/6)(s+5/6)$. And $w(\Delta\lambda)$ is not controlled by
  $w(\lambda)$ when the Picard rank exceeds one, since $w$ and the quantum
  grading are different functionals and leading terms can cancel.
\end{itemize}

Both conditions hold together iff
\[
  \boxed{\ e\,w(\Delta\lambda)<\varepsilon\ }
\]

\subsection*{The criterion is not an artifact of choosing an order}

For $C:=e\,w(\Delta\lambda)/\varepsilon\ge1$ the obstruction is not about well
ordering at all. The $z^n$-coefficient of the single product $GP$ of a
pro-Laurent gauge with a splitting gauge is $\sum_mg_mp_{n+m}$, whose terms have
weight at least $m(1-C)-nC$; for $C\ge1$ infinitely many terms have bounded
weight, so that coefficient is an infinite sum in no completion, ordered or not.
Any repair that forms this product must therefore establish $C<1$. A repair that
never forms it --- see route~2 below --- is not bound by this.

\subsection*{The criterion depends on a normalization the draft may choose}

This is the practical point, and it reverses the recommendation of the first
version of this memo. The draft fixes $w(u)=1$ and $w(s_{j,\ell})=1$
independently of $L$, and sets
$w(Q^{i_*d}u^{\rho_C\cdot d})=L(H\cdot i_*d)+\rho_C\cdot d$. Hence
$\varepsilon=1$ always, while the weights of the Novikov generators grow
linearly in $L$. The criterion is therefore
\[
  e\bigl(L(H\cdot i_*d_0)+\rho_C\cdot d_0\bigr)<c_1\cdot d_0,
\]
using that the eigenvalues of $c_1\star$ are homogeneous of quantum degree one
when $Q^d$ has degree $c_1\cdot d$ --- a convention I have checked against the
draft's own low-dimensional vanishing proof, where the $Q^\beta$-component of
$c_1\star$ maps $H^i$ to $H^{i+2-2c_1\cdot\beta}$, and against its explicit
cubic block, where $K_0$ has eigenvalues $\pm6r$, $0$, $0$ with $r=(3q)^{1/2}$
and $c_1\cdot\ell=2$. Consequences:

\begin{itemize}
\item Under the draft's instruction to choose $L$ ``so large that'' separation
  holds, the criterion fails at essentially every comparison with two distinct
  exponential factors --- including the cubic endpoint.
\item Under the smallest admissible $L$ it can hold. For a projective bundle
  with $c_1(V)\cdot\ell=0$ one may take $L=1$, giving $w(Q^\ell)=1=\varepsilon$
  and $w(\Delta\lambda)=w(q^{1/2})=1/2$.
\item Independently of $L$: separation forces $w\ge\varepsilon$ on every
  generator of $J_j$, and the image of $Q^{d_0}$ is a generator, so
  $c_1\cdot d_0\ge2$ is necessary under any admissible weight. Index one fails
  outright.
\end{itemize}

So the cheapest available repair is to choose the separating weight minimally
rather than large, and to check the resulting inequality case by case. The first
version of this memo asserted the opposite, on the false ground that the ratio
was scale-invariant; rescaling $w$ does move both rates equally, but $L$ is not
a rescaling, because the bulk generators' weights are pinned at one.

\section{The iterated receiver}
\label{sec:iterated}

The criterion above is a consequence of one architectural choice, and a reader
of the second version of this memo proposed the alternative: the draft puts the
intrinsic Novikov variables of the endpoint and the new completion variables
into a single ordered group, which is what gives an eigenvalue difference
positive weight. Nothing forces that. Treat the intrinsic Novikov fraction field
as \emph{coefficients}, and complete only in the new exceptional and bulk
directions afterwards.

Two facts make the proposal concrete, and both are visible in the draft.

\begin{itemize}
\item \emph{The unbounded negative $z$-order is bulk-driven, not
  Novikov-driven.} In the recursion (4.1a) the coefficient $G_\alpha$ has no
  power of $z$ below $-|\alpha|$, where $\alpha$ indexes the \emph{bulk}
  parameter $\eta$, and the filtration depth is acquired by substituting
  $\eta\in F^1B$. A completion in the bulk directions alone therefore already
  hosts the pro-Laurent gauge; the Novikov directions are not needed for it.
\item \emph{With the Novikov directions as coefficients, the pathology
  disappears and the criterion is satisfied rather than evaded.} Eigenvalues of
  $c_1\star$ become constants of weight zero, so $\exp(-\lambda/z)$ is not a
  unit of the receiver but a genuine exponential symbol, the constant-coefficient
  hypothesis of Definition~\ref{def:solution-algebra} holds, and
  Remark~\ref{rem:constants-fails} evaporates. Since $w(\Delta\lambda)=0$, the
  inequality $e\,w(\Delta\lambda)<\varepsilon$ holds trivially, and the product
  $GP$ of Section~\ref{sec:receiver} converges coefficientwise.
\end{itemize}

The shifts that do involve the exceptional Novikov root do not force those
directions back into the filtration. The rank-one unit twist attached to
$-(c-1)\lambda_j$, the fixed divisor shift, and the ambient target's
$\tau^\circ=O(q^{-1})$ all enter Lemma~4.5 as its $a_0$ and $a_2^\circ$ terms,
which are required only to be \emph{constant in} the positive filtration, not to
lie in it. With the Novikov directions as coefficients, $\exp(-a_0/z)$ is an
honest exponential symbol rather than the object the draft must argue about
separately.

\subsection*{What the sources actually permit}

I have checked this against the pinned sources rather than left it as a
question, and the answer is positive for the comparison and negative for the
gauge.

Iritani's Theorem~5.18 states $\Psi$ as an isomorphism of
$\C[z](\!(q^{-1/s})\!)[[Q,\widetilde\tau]]$-modules, with the change of
variables defined over $\C(\!(q^{-1/(r-1)})\!)[[Q]]$. So the ring is
\emph{Laurent} in the exceptional root and \emph{formal} in the ambient Novikov
variables and the bulk. Two consequences.

\emph{The comparison survives.} Inverting the Novikov monomials and completing
in the exceptional and bulk directions is a ring homomorphism out of that ring.
$\Psi$, its inverse, and the two inverse identities are module maps over it, so
they base-change along it unchanged, and the framed operators' entries lie in
the coefficient field, which injects. Nothing in the transport step needs
Iritani's proof, only his statement, so the formality of his ring in $Q$ is not
by itself an obstacle.

\emph{The gauge does not.} The source's own reconstruction writes
$\tau(\widetilde\tau)=\tau^\circ+t(\widetilde\tau)$ and
$\varsigma_j(\widetilde\tau)=\varsigma_j^\circ+s_j(\widetilde\tau)$, with
$t(\widetilde\tau)\in H^*(X)(\!(q^{-1})\!)[[Q,\widetilde\tau]]$ and $s_j$ of the
same shape over $H^*(Z)$. The initial terms are harmless: $\tau^\circ$ and
$\varsigma_j^\circ$ are by definition the values at $Q=\widetilde\tau=0$, hence
carry no Novikov dependence at all, and the exceptional direction remains a
completion direction as the draft already has it. But $t$ and $s_j$ vanish only
at $Q=\widetilde\tau=0$, so at $\widetilde\tau=0$ they retain pure-Novikov
terms. Those have to be removed by the normalized gauge of Lemma~4.5, and that
construction converges by topological nilpotence of the bulk shift. Inverting
the Novikov directions destroys exactly that.

So the iterated receiver is not free, and the decisive question is now small and
graded rather than general:

\begin{quote}
  Is the pure-Novikov part of the coordinate shift --- $t$ and $s_j$ restricted
  to $\widetilde\tau=0$ --- confined to $H^0\oplus H^2$?
\end{quote}

If it is, the string and divisor equations absorb it into the $a_0$ and
$a_2^\circ$ terms of Lemma~4.5, which are required only to be constant in the
positive filtration, and the iterated receiver goes through for the ambient
summand with $w(\Delta\lambda)=0$. If it is not, the Novikov directions must
stay inside the filtration, eigenvalue differences reacquire positive weight,
and the criterion of Section~\ref{sec:receiver} returns.

The grading answers this, and the answer is negative in general. Iritani's
Theorem~5.18(5) states that $\tau(\widetilde\tau)$ and
$\varsigma_j(\widetilde\tau)$ are homogeneous of degree $2$, and his conventions
are $\deg Q^d=2c_1(X)\cdot d$ and $\deg\tau_i=2-\deg\phi_i$. A pure-Novikov term
$Q^d\phi$ therefore satisfies
\[
  \deg\phi=2-2\,c_1(X)\cdot d,
\]
so $c_1\cdot d=1$ contributes in $H^0$, $c_1\cdot d=0$ in $H^2$, and
$c_1\cdot d=-1,-2,\dots$ in $H^4,H^6,\dots$; classes with $c_1\cdot d\ge2$
cannot contribute at all. Hence the required confinement to $H^0\oplus H^2$
follows exactly when every contributing effective class satisfies
$c_1\cdot d\ge0$, and it is not automatic: Iritani's theorem is stated for an
arbitrary smooth projective $X$, and after several blowups in a weak
factorization there is no reason for the intermediate ambient variety to have
nef anticanonical class.

So the iterated receiver does not by itself remove arbitrary intermediate
fourfolds from the problem. It removes them under a semipositivity hypothesis on
the ambient, which the weak factorization does not supply. This is worth stating
plainly, because it is the second proposed repair that fails on inspection, and
it fails for a reason the draft would have inherited silently.

For the center the same analysis applies, with the additional and unchanged
difficulty that its Novikov map is the noninjective one, which is why divisor
tagging exists.

The scope consequence is the reason to try this first. If the ambient endpoint
is handled this way, no inequality is needed for arbitrary intermediate
fourfolds --- the case that looked least tractable. What remains is the center
summands, and there the draft's own low-dimensional classification is
restrictive: for nef canonical class it produces an integral-$z$ gauge to a
regular-singular connection, so there is no exponential separation at all and
$w(\Delta\lambda)=0$ again; the remaining minimal cases are $\PP^1$, $\PP^2$,
ruled surfaces, and point blowups. Proving primitive-sixth vanishing directly
for those, in the coordinates Iritani's specialization actually produces, is a
smaller problem than a general transport theorem.

\section{The framing route, checked against the sources and closed}
\label{sec:hyzz}

The second version of this memo proposed the following route and called it the
probable repair. I state it as it was proposed, then report the three checks,
then state what survives. The checks were run against
arXiv:2411.02266 as cached, arXiv:2307.03696v4, and arXiv:2307.13555v3; the
full verdicts, locators and file hashes are in the companion report
\emph{The three framing-compatibility checks, run against Iritani's
decomposition}.

\subsection*{The route as proposed}

Everything above treats the pro-Laurent bulk gauge of Lemma~4.5 as unavoidable.
It may not be. Hinault, Yu, Zhang and Zhang work with $F$-bundles on the formal
$u$-disc, where a framing is a regular flat connection and, in a framing-flat
trivialization, the meromorphic connection takes the normal form
\[
  \nabla_{u\partial_u}=u\partial_u+u^{-1}K+(\mu D+H).
\]
Their Theorem~4.34 characterizes when two connections of that shape are
gauge-equivalent under
\[
  \Phi(u)\in\mathrm{GL}(H[[u]]),
  \qquad \Phi_{ij}(u)\in k[c_1,\dots,c_r][[u]],
\]
with $\Phi$ then unique given its value at $u=0$; the conditions are that the
leading operators $K$ be conjugate by a parameter-polynomial $\phi$, that the
two $\mu$ agree, and that the diagonal entries of $H$ agree after conjugation
modulo the parameter ideal. Their applications are uniqueness of the
decomposition for projective bundles and for blowups, complementing the
existence theorems of Iritani--Koto and Iritani.

The shape of that gauge is exactly what this memo has been missing. A gauge in
$\mathrm{GL}_n(K[[u]])$ has only nonnegative integral powers of the loop
coordinate, adjoins no root of $u$, and is therefore fixed by the turn
$u\mapsto e^{2\pi i}u$. Taking $\Gamma=0$, so that $\mathcal H=\Omega_V(\!(u)\!)$
and Levelt--Turrittin is available, Theorem~\ref{thm:transport} applies with no
further work:
\[
  \Phi\in\mathrm{GL}_n(K[[u]]),\quad
  \nabla'=\Phi\nabla\Phi^{-1}-\theta(\Phi)\Phi^{-1}
  \ \Longrightarrow\
  M_{\mathrm f,V}(\nabla')\ \text{and}\ M_{\mathrm f,V}(\nabla)\ \text{conjugate}.
\]
No Hahn receiver, no product $GP$, and no inequality
$e\,w(\Delta\lambda)<\varepsilon$, provided the decomposition of
Proposition~4.7 can be identified with an $F$-bundle isomorphism of the class
Theorem~4.34 governs. That identification, the matching of the $u=0$ datum, and
the separation of the center specialization were the three checks.

\subsection*{Check one: the identification does not exist as posed}

It fails three times over, and the first two are independent of the draft.

\emph{Domain.} Section~4.4 of Hinault--Yu--Zhang--Zhang is titled
\emph{Equivalence of $F$-bundles over a point}, and both of its applications
evaluate at the closed point $q=t=0$ of the base, the large-radius limit; the
proof of their Theorem~5.16 opens by observing that there the quantum product
reduces to the classical cup product, so that the leading operator $K_{\rm
split}$ of their (5.18) is a cup-product operator and hence nilpotent. The
draft's $\nu_6$ is not evaluated there. It is read where the eigenvalues of
$c_1\star$ separate --- at the cubic, where the block carries $r=(3q)^{1/2}$ and
the indicial roots $\pm1/6$, $\pm5/6$. At $q=0$ that separation collapses and
the exponential blocks the framed operator is built from do not exist. Their
Theorems~5.20 and 5.24 do reconstruct the isomorphism over the whole base from
its restriction at the limit point, but that is uniqueness of a map whose
existence is Iritani's and Iritani--Koto's theorem; it is not a transport of
framed monodromy between base points.

\emph{Hypotheses, for the blowup.} Section~4.4 assumes the generalized
eigenspaces of $K$ all have the same dimension, and the matrix calculus of
Definition~4.33 and Theorem~4.34 rests on the resulting splitting
$H_0^{\oplus m}\cong H$. The blowup decomposition is
$H^*(X)\oplus\bigoplus_{i=1}^{m-1}H^*(Z)[-2i]$, whose summands differ in
dimension whenever $\dim H^*(Z)\ne\dim H^*(X)$; the authors write a general
endomorphism there as a matrix with $\Phi_{1,1}\in\operatorname{End}(H^*(X))$
and $\Phi_{i,i}\in\operatorname{End}(H^*(Z))$, so the mismatch is explicit in
their own notation. Their blowup statement, Theorem~5.22, is accordingly called
analogous and carries \emph{no proof}; its uniqueness clause is conditional on
restricting the coefficients of $\Phi$ to a universal algebra, and for existence
they refer to Iritani--Koto. So the blowup half of Proposition~4.7 --- equation
(4.3), the half the headline theorem cannot do without --- has no proved
statement of theirs behind it. The projective-bundle half is genuinely covered:
there the eigenspaces are $m$ copies of $H^*(X)$ by their Lemma~5.8, matching
Section~4.4, and Theorem~5.16 is proved.

\emph{Purpose.} Even where Theorem~4.34 applies, its conclusion compares the
source at its base point with the target based at a shifted class $\Delta(a)$,
pinned by their (5.17) for projective bundles and (5.23) for blowups. That
shift is exactly what the pro-Laurent gauge of Lemma~4.5 exists to undo. The
route relocates the bulk displacement rather than removing it.

The two normalizations do agree on where the displacement sits, which is a
useful check on the draft's reading of the sources. In their limit-point
coordinates the eigenvalue $\lambda_i\in\C[c_1,\dots,c_m]$ is cohomology-valued,
so $\Delta(a)$ acquires components in $H^{\ge4}$ from the higher Chern classes,
while the $H^2$ components cancel identically --- their Lemma~5.8(2) gives
$m\lambda_i\equiv m\xi^{i-1}-c_1V$, cancelling the $c_1V$ on the right of
(5.17) --- and the $H^2$ part of $\Delta(a)$ is left free. In Iritani's
coordinates the same data reads
$\varsigma_j^\circ=-(c-1)\lambda_j+h_{Z,j}+O(q^{-1/(c-1)})$ with $\lambda_j$ a
\emph{scalar} multiple of $q^{1/(c-1)}$, so the $H^0$ part is that scalar, the
$H^2$ part is $h_{Z,j}$, and everything of degree at least four has moved into
the tail. Same content, different bookkeeping; the draft's split of
$\varsigma_j^\circ$ into unit twist, fixed divisor, and positive-filtration tail
is right. The tail is irreducible: for a trivial rank-two bundle all $c_i(V)$
vanish and their $\Delta(a)$ is purely $H^0\oplus H^2$, yet Iritani--Koto's
$\varsigma_j^\circ$ still carries its $O(q^{-1/r})$ tail, so even
$X\times\PP^1$ does not escape the bulk gauge.

\subsection*{Check two: the $u=0$ datum, which passes}

Their Theorems~5.20(2) and 5.24(2) determine the base map from its restriction
to the base point up to a multiplicative constant in each logarithmic direction;
the proof of their Proposition~4.31 shows why, since $df$ is recovered through
$\Psi(d\log p_i)=d\log f_i$ and integrating a logarithmic form leaves one
constant. That ambiguity is a character of the Novikov lattice: $q\mapsto cq$ is
$Q^d\mapsto c^{E\cdot d}Q^d$, which is the divisor substitution (4.1) of the
draft with $a_2^\circ$ a logarithmic multiple of the dual class. The draft
already proves such a substitution is a $z$-constant coefficient automorphism
and leaves $\nu_6$ unchanged, and already checks it is well defined on the image
monoid. The bundle map carries no ambiguity at all once its restriction is
fixed.

The ambiguity also never has to be tolerated, because the sources pin it.
Iritani--Koto's (5.11) gives
$\varsigma_j^\circ=r\lambda_j+[z^{-1}]\log\bigl(q^{c_1(V)/(rz)}F_j(1)\bigr)$ in
closed form and their (5.12) gives $\Phi_j^\circ$; Iritani's Section~5.8.1 fixes
$\Psi|_{Q=\widetilde\tau=0}$ from his Lemma~5.13 and the explicit $\lambda_j$.
One manuscript sentence citing those, and noting that any residual constant is a
divisor substitution already covered by (4.1), settles this point for good,
whichever repair eventually lands.

\subsection*{Check three: the center specialization, which stays separate}

It stays separate more firmly than the check anticipated. Their $A$-model
$F$-bundle is built throughout on the one-variable Novikov projection
$q^\beta\mapsto q^{\beta\cdot\omega}$ of their (2.21), and in the blowup
application the center copies are the maximal $F$-bundles of $Z$ associated to
$\sigma^*\omega_X$ --- the Novikov direction pulled back from the ambient
variety, which is the draft's specialization in a coarser form. No statement of
theirs can therefore certify that a specialized center module carries an
intrinsic framed monodromy. Divisor tagging remains the only route to intrinsic
center invariants.

One genuine gain came out of it. Their Lemma~2.24 proves that the collapsed
potential determines the full genus-zero potential, by expanding in the $H^2$
bulk directions and separating curve classes by their intersection numbers with
an $H^2$ basis, under the finiteness of their Assumption~2.22. That is the
mechanism of \texttt{lem:divisor-tagging}, in print, in a closely related
setting. It does not close that lemma --- it works at the level of the potential
over a base retaining all $H^2$ directions, whereas the draft needs the
separation at the level of a framed operator at a fixed bulk point, and the
draft's collapse is by the pair $(i_*,\rho_C\cdot)$ rather than by a single nef
class --- but it makes the shape of the argument a cited one.

\subsection*{What survives, and the corrected diagnosis}

Two things survive, and both are worth more than the route was.

\emph{The comparison isomorphisms need no receiver.} Iritani's $\Psi$ is an
isomorphism of $\C[z](\!(q^{-1/s})\!)[[Q,\widetilde\tau]]$-modules
(Theorem~5.18), and Iritani--Koto's $\Phi$ is induced by a
$\C[z,q][[Q,\widehat\tau]]$-module map and is an isomorphism over
$\C[z](\!(q^{-1/r})\!)[[Q,\widehat\tau]]$ (proof of their Theorem~5.1(6)). Both
are polynomial in $z$ with inverses in the same ring, hence lie in
$\mathrm{GL}_n(F(\!(z)\!))$ for $F$ the Laurent coefficient field, hence are
fixed by the turn. Theorem~\ref{thm:transport} applies to them with $\Gamma=0$,
where Levelt--Turrittin is available: no Hahn field, no product $GP$, no
inequality. The receiver problem was never a property of the comparison.

\emph{The pro-Laurent gauge is not the draft's choice.} This corrects the
sentence the second version of this memo built its recommendation on.
Iritani--Koto transport along the same object: their reconstruction in
Section~5.8 uses
$M=\bigoplus_je^{-\varsigma_j^\circ/z}M_B(\varsigma_j^\circ+s_j;Qq^{-c_1(V)/r})$,
described in their footnote as a fundamental solution for the summands with
respect to $(Q,s_0,\dots,s_{r-1})$ but not for $q$ and $z$, satisfying
$M=\mathrm{id}+O(z^{-1})$ and lying in
$\operatorname{End}[z^{-1}](\!(q^{-1/r})\!)[[Q,s]]$. Per bulk degree it is a
\emph{polynomial} in $z^{-1}$ --- which is the draft's own bound (4.1a), that
$G_\alpha$ has no power of $z$ below $-|\alpha|$ --- and it is unbounded only
after summing over bulk degrees. The sources meet this object and handle it not
by an ordered receiver but by Birkhoff factorization, splitting the $z$-regular
factor $\Phi$ from the $\mathrm{id}+O(z^{-1})$ factor $M'$ in
$(\Phi^\circ)^{-1}\circ M=M'\circ\Phi^{-1}$. Any repair should be measured
against that, not against a receiver.

\subsection*{The residual gap}

After the two survivals, the blocker is no longer transport across the
comparison. It is one statement:

\begin{quote}
  Let $T$ be smooth projective with quantum connection $\nabla$, and let
  $\tau^\bullet$ be a bulk parameter in the positive filtration --- Iritani's
  $\tau^\circ=q^{-1}[Z]+O(q^{-2})$ on the ambient summand, the
  $O(q^{-1/(c-1)})$ tail of $\varsigma_j^\circ$ together with $s_j$ on a center
  summand. Show that the framed formal monodromy of $\nabla|_{\tau=\tau^\bullet}$
  has the same primitive-sixth multiplicity as $\nabla|_{\tau=0}$.
\end{quote}

Everything else in Proposition~4.7 is proved or reduced to a regular gauge, and
\texttt{lem:divisor-tagging} inherits this gap and nothing else.

\section{Routes, in the order I would now try them}

The framing route has been removed from this list by the checks of
Section~\ref{sec:hyzz}. All five below attack the residual gap stated at the end
of that section.

\begin{enumerate}
\item \textbf{Birkhoff.} The identity $(\Phi^\circ)^{-1}\circ M=M'\circ\Phi^{-1}$
  of Iritani--Koto's Section~5.8, and its counterpart in Iritani's
  Section~5.8.2, already factors the bulk transport into a $z$-regular factor,
  which Theorem~\ref{thm:transport} handles, and a factor
  $M'=\mathrm{id}+O(z^{-1})$. What must be shown is that the second factor does
  not change the primitive-sixth multiplicity. This is a statement about a
  $z=\infty$-side gauge acting on a $z=0$ invariant, so it is not automatic; it
  is, however, sharply posed and of a standard shape, unlike the receiver
  inequality.
\item \textbf{Pointwise residues, avoiding transport altogether.} In
  Iritani--Koto's displayed conventions, $\nabla_a=\partial_a+z^{-1}(\phi_a\star)$
  and $\nabla_{z\partial_z}=z\partial_z-z^{-1}(E\star)+\mu$, two exact relations
  hold: commutativity and associativity of the quantum product give
  $[E\star,\phi_a\star]=0$, and flatness gives
  $\partial_a(E\star)=\phi_a\star+[\phi_a\star,\mu]$. The first says every bulk
  direction preserves the spectral decomposition of $E\star$, so a bulk
  displacement never mixes exponential blocks; the second identifies
  $\partial_a\lambda_i$ with the $i$-th eigenvalue of $\phi_a\star$, the
  canonical-coordinate relation. Together they invite comparing block residues
  at the two parameters instead of transporting between them --- the ``argue
  graded piece by graded piece'' possibility the earlier versions left open,
  now with the relations that make it concrete. The signs above follow the
  displayed connection of Iritani--Koto's Theorem~5.1 and should be rechecked
  against the draft's normalization before use.
\item \textbf{The low-dimensional centers directly}: primitive-sixth vanishing
  after Iritani's specialization, using the classification the draft already
  proves, rather than through a general transport statement.
\item \textbf{Divisor tagging on its own terms}: whether the divisor equation
  gives a more explicit transport than the general positive-bulk gauge of
  Lemma~4.5. Their Lemma~2.24 is the same mechanism in print and is worth
  citing; it does not close the lemma, for the reasons under check three.
\item \textbf{Normalize, then verify.} Replace ``choose $L$ so large'' by the
  smallest weight satisfying separation, adopt the single-copy receiver of
  Section~\ref{sec:receiver}, and verify $e\,w(\Delta\lambda)<\varepsilon$ case
  by case. This is fallback machinery now.
\end{enumerate}

The iterated receiver of Section~\ref{sec:iterated} is not on the list: it needs
a semipositivity hypothesis on the ambient variety that weak factorization does
not supply. It remains worth having where it applies. Restricting
Proposition~4.7 to centers satisfying the criterion is likewise available, at a
cost in generality that is the author's call.

One observation bears on how much any of these has to prove. In Dubrovin's
semisimple normalization the operator $\mu$, conjugated to canonical
coordinates, has zero diagonal, so a multiplicity-one block of $E\star$
contributes residue zero and no root of unity. If that carries over to this
normalization --- I flag it rather than assert it --- then $\nu_6$ is supported
entirely on blocks with repeated $E\star$-eigenvalue, and the routes above need
only control those. It is consistent with the draft's own cubic block, where
$K_0$ has eigenvalues $\pm6r,0,0$ and the primitive-sixth roots come from the
rank-two $\lambda=0$ block through $\det L_s=(s+1/6)(s+5/6)$.

A separate scope reduction stands on its own, unchanged: a direct quantum
K\"unneth computation for $X\times\PP^1$ together with a direct computation for
$\PP^4$ would leave the blowup step as the only operation formula the headline
theorem needs, with the genus-eight corollary waiting. Note that this cannot be
had by triviality of the bundle alone --- see the end of check one.

\section{Placement in the manuscript}
\label{sec:placement}

Stated by label because these insertions renumber Section~4. Items 1 to 3 and
item 5 are unconditional and can be written now; the checks of
Section~\ref{sec:hyzz} also make item 6 partly unconditional, since the
comparison isomorphisms themselves are $z$-polynomial and the transport lemma
applies to them with no receiver. Only the steps that pass through the
pro-Laurent bulk gauge remain conditional on closing the residual gap.

\begin{enumerate}
\item The ground field and solution algebra go immediately before ``Apply
  Levelt--Turrittin after a ramification $z=w^e$'', which already refers to a
  universal formal solution algebra that is never constructed.
\item Lemma~\ref{lem:constants} and Proposition~\ref{prop:framed-is-turn} go
  after \texttt{def:framed-sixth-multiplicity}, converting that definition into
  a characterization, with Remark~\ref{rem:constants-fails} stating the
  hypothesis that makes it available.
\item Theorem~\ref{thm:transport} becomes \texttt{lem:frame-transport}, after
  the paragraph beginning ``The frame matters'' and before
  \texttt{def:pro-laurent-gauge-group}.
\item In \texttt{lem:pv-base-change}, the asserted final clause and the
  corresponding sentence of its proof become a reference to
  \texttt{lem:frame-transport} --- but only once the receiver repair makes that
  lemma applicable there, since that lemma lives over the pro-Laurent receiver.
\item At (4.1c) in \texttt{lem:formal-base-shift}, one sentence identifying
  $M^{\mathrm{bulk}}$ as the transported framed automorphism, in the module
  frame, with the solution-frame statement cited.
\item \texttt{prop:framed-operations} cites the lemma at its three load-bearing
  steps and says which field plays the role of $\mathcal H$ at each. Two of
  those steps are now unconditional: the transport across $\Psi$ and across
  $\Phi$ needs only that both are $z$-polynomial with inverses in the same ring,
  which their statements give, so the field there is $F(\!(z)\!)$ over the
  Laurent coefficient field and no criterion is verified. The step that removes
  the positive-filtration bulk parameter is the one that waits.
\item A sentence in \texttt{prop:framed-operations} fixing the normalization of
  the coordinate change, citing Iritani--Koto (5.11)--(5.12) and Iritani's
  Section~5.8.1, and noting that any residual constant in a logarithmic Novikov
  direction is a divisor substitution already covered by (4.1). This is owed
  regardless of which repair lands, and is check two of
  Section~\ref{sec:hyzz}.
\item \texttt{lem:divisor-tagging} carries the same obstruction --- it embeds
  the gauge of \texttt{lem:formal-base-shift}, whose negative $z$-order is
  unbounded, together with every transported operator, into one ordered Hahn
  field --- and needs whatever repair the receiver gets.
\end{enumerate}

Mechanically, a new theorem-like environment must carry exactly one semantic
label and a matching row in the claim map consumed by the manuscript's
source-only gate, with an empty declaration list exactly when its coverage is
recorded as absent, and the coverage snapshot string in both verification
READMEs must be regenerated.

\section{Corrections to the earlier versions of this memo}
\label{sec:corrections}

\subsection*{From the second version to this one}

Both changes come from running the three checks against the sources.

\begin{enumerate}
\item The framing route of Section~\ref{sec:hyzz} was called the probable
  repair. It is not a repair at all: its governing theorem lives at the
  large-radius limit point, its blowup case is unproved and outside that
  theorem's hypotheses, and the gauge it supplies lands at the shifted base
  point the bulk gauge exists to undo.
\item The second version wrote that the unbounded negative loop order is the
  draft's own choice rather than something the comparison theorems impose. That
  is wrong, and it was the ground for the recommendation. Iritani--Koto
  transport along the same object, their fundamental solution $M$ of
  Section~5.8, with the same per-bulk-degree $z^{-1}$-polynomiality the draft
  derives in (4.1a).
\end{enumerate}

Two findings were gained in exchange, and they narrow the problem rather than
restate it: the comparison isomorphisms need no receiver, and the residual gap
is one displacement statement. Both are in Section~\ref{sec:hyzz}.

\subsection*{From the first version to the second}

For readers of the first version, these are the substantive changes.

\begin{enumerate}
\item The Constants Lemma was stated without the constant-coefficient
  hypothesis and is false without it; see Remark~\ref{rem:constants-fails}. The
  framed-operator characterization inherits the hypothesis.
\item The scissors were only implicit. The transport theorem has no verified
  instance in the application, and that now appears where the theorem does.
\item The claim that the criterion is scale-invariant, and the inference that
  the separating parameter $L$ cannot be tuned, were both wrong. Tuning $L$ down
  is now the first proposed route.
\item The cubic-endpoint bullet compared two different normalizations. Corrected
  above.
\item The finite-tensor objection to the receiver was wrong twice over: the
  pro-Laurent gauge lies in the Hahn factor alone, and the splitting gauge's
  coefficients lie in one finite extension of $H_{0,j}$. It is replaced by the
  gauge-relation argument.
\item The remark that a tensor product of two fields is not a domain was
  asserted as a general principle, which is false. The correct statement is that
  $\mathscr R_j$ need not be one.
\item The obstruction is now given in its order-free form, which answers two of
  the questions the first version left open: no completion helps, and the
  criterion is not an artifact of the ordered receiver.
\item Smaller repairs: the freeness proof (the relevant quotient group has
  torsion), the commutation of $\sigma$ with the derivation, the inverse of $M$,
  the range of the exponents, the ramification factor in the splitting rate, and
  the divisor-tagging observation, which is new.
\end{enumerate}

\section{Questions}

The first question of the second version --- whether the two connections
compared in Proposition~4.7 lie in the class Theorem~4.34 governs --- has been
answered no, three times over, in Section~\ref{sec:hyzz}. These replace it.

\begin{enumerate}
\item Does the factor $M'=\mathrm{id}+O(z^{-1})$ of the sources' Birkhoff
  factorization preserve the primitive-sixth multiplicity? This is now the
  question that decides the repair.
\item Can the framed spectrum at a positive-filtration bulk parameter be
  compared with the one at the origin blockwise, using
  $[E\star,\phi_a\star]=0$ and
  $\partial_a(E\star)=\phi_a\star+[\phi_a\star,\mu]$, without transporting?
\item Does the vanishing of the diagonal of $\mu$ in canonical coordinates carry
  over to this normalization, so that only repeated-eigenvalue blocks can carry
  $\nu_6$?
\item For the low-dimensional centers, can primitive-sixth vanishing be proved
  directly in the coordinates Iritani's specialization produces, given that the
  nef-canonical case already yields an integral-$z$ gauge to a regular-singular
  connection?
\item Failing all of these, what is the right restriction on centers, and does
  weak factorization stay inside it?
\end{enumerate}

\end{document}
